\documentclass[12pt,a4paper]{article}
\usepackage[margin=2.5cm]{geometry}
\usepackage{hyperref}
\usepackage{microtype}

\begin{document}
\pagestyle{empty}

\noindent
Samir Orujov\\
ADA University / ICTA / UNEC, Baku, Azerbaijan\\
Universit\'e Bretagne Sud, Vannes, France\\
\texttt{sorujov@ada.edu.az}

\bigskip
\noindent\today

\bigskip
\noindent
The Editors\\
\textit{Journal of Nonparametric Statistics}\\
Taylor \& Francis

\bigskip

\noindent Dear Editors,

\medskip

I submit for your consideration the manuscript \textbf{``Anti-Modes,
Distributional Gaps, and the Fisher Information Profile: Nonparametric
Estimation and Inference''}, which I believe is well-suited to the
\textit{Journal of Nonparametric Statistics}.

\medskip

\noindent\textbf{Topical fit.}
The paper lies squarely within nonparametric statistics.  Its core
contributions are: (i) a kernel plug-in estimator for the anti-mode
percentile with a fully derived asymptotic distribution ($n^{2/5}$-consistency,
bias--variance decomposition); (ii) a bootstrap test for the null of no
anti-mode, with Hadamard-differentiability-based consistency
(Fang \& Santos 2019); and (iii) a systematic characterisation of
density anti-modes via the Fisher information profile.  These contributions
are methodological and theoretical, with a simulation section following
standard GNST conventions.

\medskip

\noindent\textbf{Novelty.}
Despite extensive literature on mode estimation, bump-hunting, and
unimodality testing, density anti-modes---local minima of the density---have
received no systematic theoretical treatment.  The paper provides the first:
\begin{itemize}
  \item Rigorous definition and characterisation of anti-modes via equivalent
    conditions on the density, quantile density, and Fisher information
    profile (Theorem~2.3 and Corollary~2.4).
  \item Asymptotic normality result for a kernel anti-mode estimator,
    including optimal bandwidth and explicit asymptotic variance
    $V = \|K'\|^2 q(p^*)^5 / (c_0^3 \gamma^4)$, a pure constant.
  \item Bootstrap test for unimodality vs.\ the presence of an anti-mode,
    with local power analysis and finite-sample coverage guarantees.
\end{itemize}
Simulation evidence (12 DGPs, $R=500$ replications) shows RMSE slopes
near $-0.40$ (theoretical optimum), detection rates of 75--82\% for
moderate-to-sharp barriers at $n=500$ rising to 80--97\% at $n=5000$,
and bootstrap coverage close to or above 95\% at moderate sample sizes.
A compact empirical illustration on the Old Faithful geyser dataset
($n=272$) demonstrates the full inferential output---detection, location
estimate, 95\% bootstrap confidence interval, and Fisher information at
the anti-mode---on a canonical bimodal dataset.

\medskip

\noindent\textbf{Companion paper disclosure.}
A companion paper (Orujov 2025, \textit{Economics Letters}, under review)
applies the statistical framework developed here to income distributions
across countries, introducing the ``glass-barrier'' concept empirically.
The two papers are self-contained and address different audiences
(statistics vs.\ economics); the present paper contains all proofs and asymptotic theory,
which are stated without proof in the companion paper; the simulation
tables are entirely distinct.  I disclose the
companion paper in a remark in Section~2 and in the bibliography.

\medskip

\noindent\textbf{Other declarations.}
\begin{itemize}
  \item The manuscript is not under review elsewhere.
  \item This research received no specific grant from any funding agency
    in the public, commercial, or not-for-profit sectors.
  \item There are no competing interests to declare.
  \item Replication code for all Monte Carlo experiments is publicly
    available at \url{https://github.com/sorujov/Antimode-JNPS}.
  \item The manuscript has been prepared using the \texttt{interact.cls}
    Taylor \& Francis template and complies with GNST author guidelines
    (single-blind, $\leq 150$-word abstract, five keywords, MSC codes).
\end{itemize}

\medskip

I hope the manuscript is of interest to \textit{GNST} and its readership.
I am happy to address any editorial questions.

\bigskip

\noindent Yours sincerely,

\bigskip
\bigskip

\noindent Samir Orujov

\end{document}
